% Options for packages loaded elsewhere
\PassOptionsToPackage{unicode}{hyperref}
\PassOptionsToPackage{hyphens}{url}
\PassOptionsToPackage{dvipsnames,svgnames,x11names}{xcolor}
\documentclass[
  american,
  11pt,
  letterpaper,
]{article}
\usepackage{xcolor}
\usepackage[margin=1in]{geometry}
\usepackage{amsmath,amssymb}
\setcounter{secnumdepth}{-\maxdimen} % remove section numbering
\usepackage{iftex}
\ifPDFTeX
  \usepackage[T1]{fontenc}
  \usepackage[utf8]{inputenc}
  \usepackage{textcomp} % provide euro and other symbols
\else % if luatex or xetex
  \usepackage{fontspec}
  \defaultfontfeatures{Scale=MatchLowercase}
  \defaultfontfeatures[\rmfamily]{Ligatures=TeX,Scale=1}
\fi
\usepackage{lmodern}
\ifPDFTeX\else
  % xetex/luatex font selection
  \setmainfont[]{Times New Roman}
\fi
% Use upquote if available, for straight quotes in verbatim environments
\IfFileExists{upquote.sty}{\usepackage{upquote}}{}
\IfFileExists{microtype.sty}{% use microtype if available
  \usepackage[]{microtype}
  \UseMicrotypeSet[protrusion]{basicmath} % disable protrusion for tt fonts
}{}
\makeatletter
\@ifundefined{KOMAClassName}{% if non-KOMA class
  \IfFileExists{parskip.sty}{%
    \usepackage{parskip}
  }{% else
    \setlength{\parindent}{0pt}
    \setlength{\parskip}{6pt plus 2pt minus 1pt}}
}{% if KOMA class
  \KOMAoptions{parskip=half}}
\makeatother
\usepackage{longtable,booktabs,array}
\newcounter{none} % for unnumbered tables
\usepackage{calc} % for calculating minipage widths
% Correct order of tables after \paragraph or \subparagraph
\usepackage{etoolbox}
\makeatletter
\patchcmd\longtable{\par}{\if@noskipsec\mbox{}\fi\par}{}{}
\makeatother
% Allow footnotes in longtable head/foot
\IfFileExists{footnotehyper.sty}{\usepackage{footnotehyper}}{\usepackage{footnote}}
\makesavenoteenv{longtable}
\ifLuaTeX
\usepackage[bidi=basic,shorthands=off]{babel}
\else
\usepackage[bidi=default,shorthands=off]{babel}
\fi
\ifPDFTeX
\else
\babelfont{rm}[]{Times New Roman}
\fi
\ifLuaTeX
  \usepackage{selnolig} % disable illegal ligatures
\fi
\setlength{\emergencystretch}{3em} % prevent overfull lines
\providecommand{\tightlist}{%
  \setlength{\itemsep}{0pt}\setlength{\parskip}{0pt}}
\usepackage{fancyhdr}
\usepackage{enumitem}
\usepackage{caption}
\usepackage{bm}
\usepackage{float}
\usepackage{microtype}
\usepackage{xurl}
\usepackage{ragged2e}

\pagestyle{fancy}
\fancyhf{}
\fancyfoot[C]{\thepage}
\renewcommand{\headrulewidth}{0pt}

\setlength{\parindent}{0pt}
\setlength{\parskip}{6pt plus 2pt minus 1pt}
\setlist{nosep,leftmargin=*}
\captionsetup{font=small,labelfont=bf}
\setlength{\emergencystretch}{3em}
\sloppy
\hyphenation{im-ple-men-ta-tion op-er-a-tion-al-i-za-tion}
\usepackage{bookmark}
\IfFileExists{xurl.sty}{\usepackage{xurl}}{} % add URL line breaks if available
\urlstyle{same}
\hypersetup{
  pdftitle={Physics in Review: Quantum Entropic-Force Models, QGTCD, and Super Information Theory},
  pdfauthor={Micah Blumberg},
  pdflang={en-US},
  colorlinks=true,
  linkcolor={blue},
  filecolor={Maroon},
  citecolor={blue},
  urlcolor={blue},
  pdfcreator={LaTeX via pandoc}}

\author{}
\date{}

\begin{document}

\begin{center}

{\LARGE\bfseries Physics in Review: Quantum Entropic-Force Models, QGTCD, and Super Information Theory\par}

\vspace{0.35em}

{\large\itshape A four-lens, atomized priority comparison with longitudinal source control\par}

\vspace{0.8em}

{\small\textbf{Author: Micah Blumberg}\par}

{\small Self Aware Networks Research Institute\par}

{\small\url{https://selfawarenetworks.com}\par}

{\small November 20, 2026 (Version 2 Status Update, Conditional Thermal Calibration, and Regime Analysis to DOI 10.5281/zenodo.21693943)\par}

\end{center}

\vspace{0.65em}

\begin{abstract}

Carney, Karydas, Scharnhorst, Singh, and Taylor construct local and nonlocal microscopic quantum models in which position-dependent mediator free energy yields a Newtonian interaction, while mediator-bath dynamics also generates force noise, heating, decoherence, and constraints on gravitational-entanglement witnesses. This paper compares that 2025 quantum entropic-force program with Micah Blumberg's earlier Quantum Gradient Time Crystal Dilation (QGTCD) genealogy and the first public version of Super Information Theory (SIT). The comparison is not all-or-nothing. It applies four separate lenses: bounded mechanism equivalence; narrative and causal-interface architecture; conclusion-first forensic reconstruction; and longitudinal comparison of the target program before, during, and after Blumberg's publications.

The dated record supports a specific, intent-neutral priority result. Blumberg publicly described in July 2022 a microscopic mass / time-density proposal in which a spatial difference in temporal opportunity changes otherwise symmetric directional probabilities and produces motion toward mass. Three 2024 SVGN articles restated the chain as locally unequal time measure, biased random motion, and an emergent gravitational trajectory. Carney et al.~submitted their quantum entropic-force paper on February 24, 2025. Its internal mechanism differs, but the difference does not erase the bounded shared architecture: a mass-associated microscopic statistical state produces a spatial scalar gradient, the gradient biases or directs motion, and the macroscopic output is Newtonian attraction without requiring fundamental graviton exchange at the compared interface. The atom ledger records six Micah-earlier partial-priority nodes in that shared causal chain. It separately preserves Carney et al.'s explicit Hamiltonians, mediator-bath reduction, Lindblad operators, quantitative noise bounds, and entanglement-witness analysis as target-specific advances.

A prior private comparison generated lists of fourteen and sixteen candidate correspondences. Those lists are useful discovery material, but their proposed multiplication of uniform match probabilities is invalid because the atoms are dependent, the null sampling frame is unspecified, and the relevant literature already contains named thermodynamic-gravity priors. This reconstruction therefore reports the probability layer as \textbf{quantitative-not-calibrated}. The strongest defensible conclusion is source-audited priority on the bounded six-node QGTCD causal sequence, plus a longitudinal shift in the Carney / Taylor gravity-test lineage from assumed Newtonian or graviton interactions in 2021-2024 to an explicit microscopic statistical-gradient account of Newtonian attraction in 2025. The newest-source check, repeated on July 27, 2026, found that the entropic-force paper still has one arXiv version, while the related April 2026 interferometer revision remains a parallel graviton-effective-field-theory strand. No claim about access, copying, awareness, or intent is made.

\end{abstract}

\noindent\textbf{Keywords:} entropic gravity; quantum gravity; QGTCD; Super Information Theory; free-energy gradient; temporal opportunity; Lindblad dynamics; priority; theory comparison; longitudinal signature migration

\section{1. Scope and claim contract}\label{scope-and-claim-contract}

Two theories can be written in different mathematical languages yet perform one of the same bounded operations. They can also share a final prediction while reaching it through genuinely different internal machinery. A useful comparison must therefore solve two problems at once: it must not miss a shared causal subchain merely because the labels differ, and it must not inflate that subchain into identity of the complete theories.

The established scientific practice is to compare assumptions, state variables, transformations, limits, and observables. This paper first explains those components in ordinary language. It then compares what each program receives as input, what changes internally, and what output follows. Only after that operation is clear does the paper use the compact labels \textbf{bounded mechanism equivalence}, \textbf{causal-interface architecture}, \textbf{conclusion-first reconstruction}, and \textbf{longitudinal signature migration}. The labels summarize an already explained comparison; they do not substitute for one.

This paper makes four kinds of statement and keeps them separate.

\begin{enumerate}
\def\labelenumi{\arabic{enumi}.}
\tightlist
\item
  \textbf{Documentary claims} identify what a dated source says and when a fixed public version existed.
\item
  \textbf{Comparison claims} map bounded inputs, operations, and outputs across theories without requiring identical terminology or ontology.
\item
  \textbf{Scientific claims} describe what a model proves under its assumptions and what remains unproved in nature.
\item
  \textbf{Historical-causal claims} would concern access, influence, copying, or intent. This paper makes none of these.
\end{enumerate}

The central historical claim is narrow but important: Blumberg's public QGTCD genealogy contains a six-node partial causal sequence before the 2025 quantum entropic-force paper. The target paper adds substantial new mathematical machinery. Added machinery is a residual, not an automatic cancellation of the shared sequence.

The comparison is governed by a patent-atom discipline: every conclusion is decomposed into inputs, transformations, outputs, and constraints. A full-operator mismatch cannot turn every shared sub-operation into zero. Conversely, a shared output cannot erase a genuinely different internal transform. The result is therefore not ``the same theory'' or ``no relation.'' It is an atomized ledger of what matches, what predates, what is supplied by named earlier literature, what remains target-specific, and what remains unresolved.

\section{2. Primary source chronology}\label{primary-source-chronology}

\subsection{2.1 Blumberg's dated sequence}\label{blumbergs-dated-sequence}

The present audit distinguishes the beginning of Blumberg's physics branch from the first recovered probability mechanism.

\begin{itemize}
\tightlist
\item
  The governing genealogy map was checked for the earlier SAN and Neo Mind Cycle ancestry. No 2011, 2012, or 2017 source is used here to backdate the QGTCD gravity operator. The paper-relevant physics branch begins in the recovered 2021 records, and the explicit directional-probability mechanism begins later {[}9{]}.
\item
  On November 11, 2021, \texttt{a0033z.md} recorded an authored attempt at an ``equation for quantum gravity.'' This dates the branch, but it does not contain the later normalized QGTCD kernel {[}9{]}.
\item
  On June 16, 2022, a public Git revision stated that gravity is gradient time dilation at particle and atomic scale {[}9{]}.
\item
  On July 26, 2022, \texttt{a0253z.md} described mass as time-crystal-like, associated mass with a time- or mass-density gradient, and described motion toward the region of greater density {[}10{]}.
\item
  On July 30, 2022, a public Git revision made the stochastic operation explicit: otherwise random directional motion receives greater odds toward mass because additional temporal measure is treated as additional available spacetime measure {[}11{]}.
\item
  On November 28, 2022, the name Quantum Gradient Time Crystal Dilation was fixed in the public repository genealogy {[}9{]}.
\item
  On January 29, February 15, and March 28, 2024, public SVGN articles restated the mechanism as quantized local regions, unequal time measure, changed trajectory odds, and a many-step gravitational path {[}12-14{]}.
\item
  On March 30, 2024, \nolinkurl{QGTCDpostulates.md} gave the compact two-part architecture: symmetric deviations cancel in a uniform environment; a time-density gradient biases the mean trajectory {[}9{]}.
\item
  On February 10, 2025, SIT version 1 was publicly posted on Figshare. It described information and coherence as active physical structure and linked information-density gradients to time dilation and gravitational effects {[}15{]}. This source is later than QGTCD and only fourteen days before the target arXiv submission; it is not used to backdate the 2022 QGTCD operation.
\item
  In 2026, a separate QGTCD paper supplied one mathematically explicit endpoint-weighted Markov realization and its drift-diffusion generator. That formalization is later and must not be projected backward into the 2022 prose {[}16{]}.
\end{itemize}

The earliest public mechanism relevant to this comparison is therefore July 2022, not the November 2021 branch marker and not the 2026 normalized kernel.

\subsection{2.2 The Carney / Taylor target sequence}\label{the-carney-taylor-target-sequence}

The target lineage is not a single static paper.

{\def\LTcaptype{none} % do not increment counter
\begingroup\footnotesize\setlength{\tabcolsep}{4pt}
\begin{longtable}[]{@{}
  >{\RaggedRight\arraybackslash}p{(\linewidth - 6\tabcolsep) * \real{0.2500}}
  >{\RaggedRight\arraybackslash}p{(\linewidth - 6\tabcolsep) * \real{0.2500}}
  >{\RaggedRight\arraybackslash}p{(\linewidth - 6\tabcolsep) * \real{0.2500}}
  >{\RaggedRight\arraybackslash}p{(\linewidth - 6\tabcolsep) * \real{0.2500}}@{}}
\toprule\noalign{}
\begin{minipage}[b]{\linewidth}\RaggedRight
Date
\end{minipage} & \begin{minipage}[b]{\linewidth}\RaggedRight
Checked target source
\end{minipage} & \begin{minipage}[b]{\linewidth}\RaggedRight
Dominant mechanism in the checked source
\end{minipage} & \begin{minipage}[b]{\linewidth}\RaggedRight
Entropic free-energy gravity present?
\end{minipage} \\
\midrule\noalign{}
\endhead
\bottomrule\noalign{}
\endlastfoot
2021-08-18 & Carney, Muller, and Taylor, atom-interferometer proposal {[}2{]} & Assumes the Newtonian potential operator and asks whether gravity generates entanglement; develops collapse-and-revival sensing with a mechanical oscillator & No \\
2023-08-24 & Carney, Domcke, and Rodd, graviton detection {[}3{]} & Tests what graviton clicks or gravitational-wave measurements can establish about quantization & No \\
2024-09-05 & Carney, Karydas, and Sivaramakrishnan, interferometer response v1 {[}4{]} & Calculates the vacuum response of low-energy graviton effective field theory & No \\
2024-12-06 & Carney and Matsumura, classical-quantum scattering {[}5{]} & Tests relativistic quantum matter coupled to a classical stochastic mediator and finds simple gravitational versions observationally inadequate & No free-energy entropic construction \\
2025-02-24 & Carney et al., quantum mechanics of entropic forces v1 {[}1{]} & Mass positions alter mediator spectra; mediator free energy yields a Newtonian interaction; open dynamics yields noise and decoherence & Yes \\
2026-04-12 & Interferometer-response v2 {[}4{]} & Published-version revision of the 2024 graviton-EFT calculation; basic conclusion unchanged & No; parallel EFT strand remains \\
\end{longtable}
\endgroup
}

This sequence supplies a bounded longitudinal result. The 2021-2024 checked stages contain oscillators, decoherence, gravity tests, gravitons, and classical-quantum alternatives, but not the target paper's causal chain from position-dependent microscopic mediator statistics through free-energy extremization to Newtonian attraction. That chain appears in the February 2025 target after Blumberg's 2022 and 2024 public QGTCD sources. Coauthor continuity is strongest through Daniel Carney and Jacob Taylor, but the author set changes across papers; the longitudinal claim therefore concerns a \textbf{Carney / Taylor gravity-test program}, not an assertion that one unchanged five-author team migrated as a unit.

The newest source check also prevents a misleading claim. The 2026 revision of the interferometer paper remains a graviton-EFT calculation. It does not replace the 2025 entropic paper, and it shows that the research program retains parallel strands rather than moving wholly into one ontology.

As of July 27, 2026, the target arXiv record still contains only the February 24, 2025 version and links that version to the published \emph{Physical Review X} article {[}1{]}. The April 12, 2026 version of the interferometer paper explicitly describes itself as the published version with presentation edits and unchanged basic conclusions {[}4{]}. This is a newest-checked-source result for the bounded Carney / Karydas strand, not a claim that every later publication by every coauthor has been exhausted.

\section{3. What the quantum entropic-force paper establishes}\label{what-the-quantum-entropic-force-paper-establishes}

\subsection{3.1 Microscopic free-energy construction}\label{microscopic-free-energy-construction}

Carney et al.~divide the model into massive bodies, microscopic mediators, and a bath. The mediator Hamiltonian depends on the positions of the massive bodies. For a thermal mediator at temperature \(T\), the Helmholtz free energy is

\begin{equation}
\mathcal{A}(\mathbf{x},T)=U(\mathbf{x},T)-T S(\mathbf{x},T),
\tag{1}
\end{equation}

and the effective force is

\begin{equation}
\mathbf{F}_{\mathrm{QEM}}(\mathbf{x})=-\nabla_{\mathbf{x}}\mathcal{A}(\mathbf{x},T).
\tag{2}
\end{equation}

The mediator spectra and couplings are selected so that the coherent effective term contains the Newtonian interaction in the model's controlled regime {[}1{]}. This is a genuine constructive result. It shows that explicit quantum-statistical mediators can realize Newtonian attraction. It does not show that all informational or entropic microphysics must yield gravity, nor that nature uses these mediators.

\subsection{3.2 Local and nonlocal realizations}\label{local-and-nonlocal-realizations}

The nonlocal model associates mediator degrees of freedom with mass pairs. The local model distributes qubits or analogous degrees of freedom through space, allowing a mass to change a surrounding mediator state. These are not merely two vocabularies for one already-proved field. They have different locality structures, parameters, constraints, and experimental consequences {[}1{]}.

\subsection{3.3 Open-system reduction}\label{open-system-reduction}

The mediators must remain near their prescribed statistical state while the masses move. Carney et al.~therefore introduce a bath and a hierarchy of timescales. Under Born-Markov and adiabatic-elimination assumptions, tracing out bath and mediator degrees of freedom produces a reduced mass evolution of the schematic form

\begin{equation}
\dot\rho_M=-i[H_M+V_N,\rho_M]
+\sum_j \mathcal{D}[K_j]\rho_M,
\tag{3}
\end{equation}

where

\begin{equation}
\mathcal{D}[K]\rho=K\rho K^\dagger-
\frac{1}{2}\{K^\dagger K,\rho\}.
\tag{4}
\end{equation}

The coherent term supplies the effective interaction. The dissipative terms supply force fluctuations, heating, and decoherence. Lindblad form itself is not unique to entropic gravity; the target paper's evidential contribution lies in the specific jump operators, scales, constraints, and observables derived from its microscopic assumptions.

\subsection{3.4 Experimental discrimination}\label{experimental-discrimination}

The target paper identifies parameter-dependent departures from inverse-square behavior, force-noise spectra, heating, spatial decoherence, and entanglement-witness behavior. These predictions allow observations to exclude model regions. They do not constitute a detection of entropic gravity. They are nevertheless more explicit than the historical 2022 QGTCD prose and must remain credited as target-specific advances.

\section{4. The bounded QGTCD mechanism}\label{the-bounded-qgtcd-mechanism}

The 2022-2024 QGTCD source sequence can be stated without importing the later 2026 mathematics:

\[
\begin{aligned}
\text{mass-associated temporal structure}
&\longrightarrow
\text{spatially unequal temporal opportunity}\\
&\longrightarrow
\text{unequal directional probability}
\longrightarrow
\text{mean motion toward mass}.
\end{aligned}
\tag{5}
\]

The 2026 formalization chooses a positive temporal-opportunity field \(\tau\), defines \(u=\log(\tau/\tau_0)\), and tilts a symmetric nearest-neighbor walk by endpoint weight. One explicit kernel is

\begin{equation}
P_\beta(x,y)=
\frac{\mathbf{1}_{y\in N(x)}\exp[\beta u(y)]}
{\sum_{z\in N(x)}\exp[\beta u(z)]}.
\tag{6}
\end{equation}

Under a regular-lattice diffusive scaling, its leading generator is

\begin{equation}
\begin{aligned}
& L f=D\nabla^2 f+2\beta D\,\nabla u\cdot\nabla f, \\
& b_{\mathrm{QGTCD}}=2\beta D\nabla u.
\end{aligned}
\tag{7}
\end{equation}

Equations (6)-(7) are a 2026 model choice, not equations quoted from 2022. They clarify the historical operator and expose its obligations: it is classical, overdamped, and dependent on a diffusion scale unless further dynamics restore universality {[}16{]}.

\section{5. Lens 1: bounded mechanism equivalence}\label{lens-1-bounded-mechanism-equivalence}

Mechanism comparison does not require that \(\mathcal{A}\) and \(u\) be the same physical object. It asks whether a declared translation preserves a bounded operation and an explicitly typed output.

At the selected scalar-gradient interface, define the black-box operator

\begin{equation}
\mathsf{G}[s;\lambda]=-\lambda\nabla s.
\tag{8}
\end{equation}

QEM instantiates this operator with \(s=\mathcal{A}\) and a force-valued output. The QGTCD drift model instantiates it with \(s=-u\), \(\lambda=2\beta D\), and a velocity-valued output. Because force and drift have different physical units, a valid comparison must include both a scalar translation and an output-unit map. One such declared interface map is

\begin{equation}
\begin{aligned}
& \mathcal{T}_s:\mathcal{A}\mapsto-c_Au, \\
& \mathcal{T}_o:\mathbf{F}\mapsto\kappa\mathbf{F}, \\
& \kappa=\frac{2\beta D}{c_A},
\quad c_A>0.
\end{aligned}
\tag{9}
\end{equation}

Under this calibration, \(\mathcal{T}_o[-\nabla\mathcal{A}]=2\beta D\nabla u\). The direction and scalar-gradient form of the mean response commute across the interface. This is a real functional correspondence: both models convert a spatial scalar or statistical imbalance into directed motion.

The correspondence stops at that interface. QEM's \(\mathcal{A}\) is a thermodynamic free energy derived from a quantum mediator state. QGTCD's \(u\) is the logarithm of a temporal-opportunity field in a classical Markov model. QEM yields a Hamiltonian term plus Lindblad residuals. QGTCD yields a drift-diffusion generator. The constants \(c_A\) and \(\kappa\) are declared comparison calibrations, not quantities derived by either theory from the other. Equation (9) is therefore an interface translation, not an identity of state spaces, dimensions, quantization, covariance, or noise.

The correct Lens-1 disposition is \textbf{bounded partial pass}. Draft 6's conclusion of mathematical non-equivalence is retained for the full models, but it cannot erase the shared scalar-gradient operation.

\section{6. Lens 2: narrative and causal-interface architecture}\label{lens-2-narrative-and-causal-interface-architecture}

The two programs can be black-boxed as causal interface sequences.

\subsection{6.1 QGTCD sequence}\label{qgtcd-sequence}

{\def\LTcaptype{none} % do not increment counter
\begingroup\footnotesize\setlength{\tabcolsep}{4pt}
\begin{longtable}[]{@{}
  >{\RaggedRight\arraybackslash}p{(\linewidth - 2\tabcolsep) * \real{0.5000}}
  >{\RaggedRight\arraybackslash}p{(\linewidth - 2\tabcolsep) * \real{0.5000}}@{}}
\toprule\noalign{}
\begin{minipage}[b]{\linewidth}\RaggedRight
Interface role
\end{minipage} & \begin{minipage}[b]{\linewidth}\RaggedRight
QGTCD content
\end{minipage} \\
\midrule\noalign{}
\endhead
\bottomrule\noalign{}
\endlastfoot
Input & Mass-associated oscillatory / time-crystal-like structure and a spatial temporal-density difference \\
Transform 1 & Unequal temporal opportunity changes the relative availability of neighboring directions \\
Transform 2 & Repeated local probability differences accumulate into a mean trajectory \\
Output & Directed motion toward mass; proposed gravitational attraction \\
Residual obligations & Quantum derivation, inertial dynamics, universality, covariance, light propagation, empirical discriminator \\
\end{longtable}
\endgroup
}

\subsection{6.2 QEM sequence}\label{qem-sequence}

{\def\LTcaptype{none} % do not increment counter
\begingroup\footnotesize\setlength{\tabcolsep}{4pt}
\begin{longtable}[]{@{}
  >{\RaggedRight\arraybackslash}p{(\linewidth - 2\tabcolsep) * \real{0.5000}}
  >{\RaggedRight\arraybackslash}p{(\linewidth - 2\tabcolsep) * \real{0.5000}}@{}}
\toprule\noalign{}
\begin{minipage}[b]{\linewidth}\RaggedRight
Interface role
\end{minipage} & \begin{minipage}[b]{\linewidth}\RaggedRight
QEM content
\end{minipage} \\
\midrule\noalign{}
\endhead
\bottomrule\noalign{}
\endlastfoot
Input & Massive-body positions coupled to qubit or oscillator mediator degrees of freedom \\
Transform 1 & Positions change mediator spectra, energy, entropy, and free energy \\
Transform 2 & Free-energy extremization creates a coherent Newtonian term; bath reduction creates dissipative terms \\
Output & Newtonian attraction plus force noise, heating, decoherence, and witness constraints \\
Residual obligations & Relativistic completion, model selection, physical mediator identification, surviving parameter region \\
\end{longtable}
\endgroup
}

The common narrative architecture is

\[
\begin{aligned}
\text{mass configuration}
&\rightarrow
\text{microscopic statistical state}\\
&\rightarrow
\text{spatial scalar imbalance}
\rightarrow
\text{directed Newtonian behavior}.
\end{aligned}
\tag{10}
\]

That four-stage sequence passes the narrative lens. The target's bath, Lindblad terms, and explicit mediator spectra enlarge the sequence; they do not reverse its shared conclusion. At the same time, the additional target machinery is not decorative. It supplies predictions the earlier QGTCD sources did not contain.

\section{7. Lens 3: conclusion-first forensic reconstruction}\label{lens-3-conclusion-first-forensic-reconstruction}

The conclusion-first lens starts with the endpoint and works backward.

\begin{enumerate}
\def\labelenumi{\arabic{enumi}.}
\tightlist
\item
  \textbf{Matching conclusion:} both programs seek Newtonian attraction from a deeper statistical or informational substrate rather than treating fundamental virtual-graviton exchange as the only available explanation at the compared nonrelativistic interface.
\item
  \textbf{First backward step:} both place a spatial gradient between the microscopic substrate and directed motion.
\item
  \textbf{Second backward step:} both make mass configuration alter the relevant microscopic or statistical state.
\item
  \textbf{Third backward step:} both use many local or microscopic degrees of freedom to produce the macroscopic response.
\item
  \textbf{Residual test:} QEM's mediator equilibrium, heat bath, adiabatic elimination, and Lindblad operators are absent from historical QGTCD; QGTCD's temporal-opportunity interpretation and directional-odds rule are absent from QEM.
\end{enumerate}

Changing the name of the scalar, the ontology of the mediator, or the mathematical packaging cannot by itself erase steps 1-4. Equally, the shared steps do not license a claim that step 5 is identical. The Lens-3 disposition is therefore \textbf{bounded partial pass with explicit two-way residuals}.

No public evidence reviewed here establishes that the target authors saw Blumberg's work. The forensic lens is used to prevent a false negative in content comparison, not to infer access or intent.

\section{8. Lens 4: longitudinal signature migration}\label{lens-4-longitudinal-signature-migration}

The longitudinal question is whether earlier checked work in a stable target program lacks a bounded signature, Blumberg's dated work occupies the middle, and later target work contains it.

The 2021 Carney-Muller-Taylor paper already contains a mechanical oscillator, thermal-state discussion, decoherence / recoherence, and an assumed Newtonian gravitational operator {[}2{]}. It does \textbf{not} derive gravity from the oscillator's free energy. The 2023 graviton-detection paper and 2024 interferometer-vacuum paper remain focused on graviton or effective-field observables {[}3,4{]}. The 2024 classical-quantum-scattering paper studies a different stochastic-mediator route and rejects a simple gravitational implementation observationally {[}5{]}.

Blumberg's July 2022 and 2024 QGTCD sources occupy the middle with the chain in Equation (5). The February 2025 QEM paper then introduces the later target chain in Equation (10) using qubits or oscillators, a spatially varying statistical scalar, and a Newtonian output {[}1{]}.

The Lens-4 result is a \textbf{bounded program-level pass} for the appearance of the microstate-to-statistical-gradient-to-Newtonian sequence after Blumberg's public QGTCD sources. It is not a pass for the exact QGTCD temporal-opportunity operator, which remains absent from QEM. It is not evidence of historical following. It is a dated content trajectory.

\section{9. Atom-by-atom genealogy and priority ledger}\label{atom-by-atom-genealogy-and-priority-ledger}

The following table is the controlling comparison. ``Micah earlier'' means earlier than the February 24, 2025 QEM arXiv submission on the bounded atom stated in that row. It does not mean earlier than every broad precursor in physics.

{\def\LTcaptype{none} % do not increment counter
\begingroup\footnotesize\setlength{\tabcolsep}{4pt}
\begin{longtable}[]{@{}
  >{\RaggedRight\arraybackslash}p{(\linewidth - 8\tabcolsep) * \real{0.2000}}
  >{\RaggedRight\arraybackslash}p{(\linewidth - 8\tabcolsep) * \real{0.2000}}
  >{\RaggedRight\arraybackslash}p{(\linewidth - 8\tabcolsep) * \real{0.2000}}
  >{\RaggedRight\arraybackslash}p{(\linewidth - 8\tabcolsep) * \real{0.2000}}
  >{\RaggedRight\arraybackslash}p{(\linewidth - 8\tabcolsep) * \real{0.2000}}@{}}
\toprule\noalign{}
\begin{minipage}[b]{\linewidth}\RaggedRight
ID
\end{minipage} & \begin{minipage}[b]{\linewidth}\RaggedRight
Bounded atom
\end{minipage} & \begin{minipage}[b]{\linewidth}\RaggedRight
Blumberg source
\end{minipage} & \begin{minipage}[b]{\linewidth}\RaggedRight
QEM source
\end{minipage} & \begin{minipage}[b]{\linewidth}\RaggedRight
Disposition
\end{minipage} \\
\midrule\noalign{}
\endhead
\bottomrule\noalign{}
\endlastfoot
A01 & Microscopic oscillatory or time-varying substrate associated with gravitational behavior & 2022 mass / time-crystal-like QGTCD proposal {[}10,\allowbreak{}11{]} & Qubits or oscillators as gravity-generating mediators {[}1{]} & \textbf{Micah-earlier partial}; broad oscillator physics is named prior art \\
A02 & Mass configuration produces a spatial change in a microscopic / statistical state & 2022 mass / time-density gradient; 2024 local quantized-region restatements {[}10-14{]} & Mass positions change mediator frequencies and free energy {[}1{]} & \textbf{Micah-earlier partial} \\
A03 & A spatial scalar / statistical gradient controls directed motion & 2022 greater temporal measure changes directional odds {[}11{]} & Effective force is \(-\nabla\mathcal{A}\) {[}1{]} & \textbf{Micah-earlier partial}, mechanism Lens-1 pass \\
A04 & Repeated microscopic bias yields motion toward mass and a Newtonian-like macroscopic output & 2022 motion toward mass; 2024 random-walk / gravitational-path explanation {[}10-14{]} & Newtonian interaction from mediator free energy {[}1{]} & \textbf{Micah-earlier conceptual output}; target earlier on the explicit quantum derivation \\
A05 & Gravity may be an emergent statistical response rather than only fundamental virtual-quanta exchange & 2022-2024 probability-distortion account {[}10-14{]} & Entropic rearrangement rather than virtual quanta in the compared construction {[}1{]} & \textbf{Micah-earlier partial on his specific route}; Jacobson / Verlinde and other entropic-gravity work are earlier broad priors {[}6,\allowbreak{}7{]} \\
A06 & A local discretized spatial realization can accumulate into macroscopic gravitational behavior & 2024 pixel / quadrant / local random-walk explanation {[}12-14{]} & Local spatial lattice of mediator qubits {[}1{]} & \textbf{Micah-earlier partial}; exact lattice Hamiltonian remains target-specific \\
A07 & Information / entropy is used to frame emergent gravity & QGTCD leads into SIT; SIT v1 public 2025-02-10 {[}15{]} & Entropic rearrangement of information, 2025-02-24 {[}1{]} & Shared, but \textbf{named-prior constrained} by Jacobson and Verlinde; no exclusive priority assigned \\
A08 & Thermal bath and adiabatic separation maintain mediator state & No 2022 QGTCD equivalent & Explicit bath and timescale hierarchy {[}1{]} & \textbf{Target-specific advance} \\
A09 & Microscopic reduction yields explicit Lindblad force noise, heating, and decoherence & Historical QGTCD has stochasticity but no quantum jump operators & Explicit model-derived dissipators {[}1{]} & \textbf{Target-specific advance}; generic open-system gravity methods are earlier field tools {[}8{]} \\
A10 & Both local and nonlocal quantum mediator constructions are worked out & Historical QGTCD supplies no paired quantum constructions & Explicit local and nonlocal QEM models {[}1{]} & \textbf{Target-specific advance} \\
A11 & Entanglement witnesses delimit gravitational-model parameter space & Not in the 2022 QGTCD operator & Detailed witness analysis building on earlier proposals {[}1,\allowbreak{}17,\allowbreak{}18{]} & \textbf{Target-specific analysis with named priors} \\
A12 & Existing and near-term experiments exclude or test parameter regions quantitatively & Historical QGTCD is qualitative; current tests are later & Noise, heating, inverse-square, decoherence, and witness constraints {[}1{]} & \textbf{Target-specific advance} \\
A13 & Temporal opportunity itself is the latent state and changes directional probabilities & Core 2022-2024 QGTCD delta {[}10-14{]} & Absent; QEM uses thermodynamic free energy & \textbf{Micah-only residual} \\
A14 & Coherence / time-density field architecture extends the gravity account into SIT & SIT v1 and later versions {[}15{]} & Absent as a field triple; QEM has a different mediator ontology & \textbf{Micah-only residual} \\
\end{longtable}
\endgroup
}

The table contains six Micah-earlier partial-priority rows: A01-A06. These are not six independent statistical trials. A02, A03, and A04 are causally linked; A01 and A06 partly specify the substrate of the same chain. The count is a content-coverage count, not a probability calculation.

\section{10. Category-theoretic and process translation}\label{category-theoretic-and-process-translation}

The comparison can be made precise enough to show both the pass and its boundary.

Let a QEM process expose mass configuration \(X\) as input, hide mediator and bath states \(M,B\), and expose mean force and covariance as output:

\begin{equation}
\mathsf{Q}: X
\xrightarrow{\text{couple}}
(X,M,B)
\xrightarrow{\text{thermalize / trace}}
(F_Q,\Sigma_Q).
\tag{11}
\end{equation}

Let a QGTCD process expose position and temporal field \((X,u)\), hide finite-step transitions, and expose mean drift and covariance:

\begin{equation}
\mathsf{M}: (X,u)
\xrightarrow{P_u}
\text{path ensemble}
\xrightarrow{\text{coarse grain}}
(b_M,\Sigma_M).
\tag{12}
\end{equation}

Define a typed black-box projection \(\mathcal{B}\) that forgets internal ontology and retains only the declared scalar field, mean gradient response, and second-order fluctuation surface:

\begin{equation}
\begin{aligned}
& \mathcal{B}(\mathsf{Q})=(\mathcal{A},-\nabla\mathcal{A},\Sigma_Q), \\
& \mathcal{B}(\mathsf{M})=(u,2\beta D\nabla u,\Sigma_M).
\end{aligned}
\tag{13}
\end{equation}

After the scalar and output translations in Equation (9), the mean-gradient component can commute. The covariance component does not generally commute: QEM's \(\Sigma_Q\) comes from quantum mediator-bath dynamics, while QGTCD's \(\Sigma_M\) comes from a classical transition kernel. In category-theoretic language, \(\mathcal{B}\) is only a proposed partial functor on the explicitly exposed interface; the source and target categories for the complete models have not been shown equivalent. Therefore:

\begin{itemize}
\tightlist
\item
  the \textbf{mean-gradient subdiagram passes} after explicit calibration;
\item
  the \textbf{full stochastic / quantum diagram does not yet pass};
\item
  no full natural isomorphism, weak bisimulation, or empirical equivalence has been proved.
\end{itemize}

This is stronger than a vocabulary comparison and more honest than declaring full mathematical identity. It identifies the exact morphism preserved by the translation and the exact output component that remains unmatched.

\section{11. Probability and derivation assessment}\label{probability-and-derivation-assessment}

\subsection{11.1 Why the old numerical claim fails}\label{why-the-old-numerical-claim-fails}

A private 2025 comparison identified sixteen candidate equivalences and proposed a calculation of \(5^{16}\) by assigning every item an independent null probability of \(0.2\). A second draft listed fourteen interface correspondences. The source-recovery value of those drafts is substantial: they preserve translations, candidate atoms, and an earlier recognition of the shared causal architecture. Their numerical inference is not valid.

The items are not independent. ``Free-energy gradient,'' ``Newtonian scalar potential,'' ``directed force,'' and ``Newtonian output'' are parts of one causal chain. ``Qubit / oscillator mediators,'' ``many-body substrate,'' and ``local lattice'' also overlap. The null sampling frame was not specified, and thermodynamic or informational gravity already has named prior literature. Multiplying a chosen probability across these dependent descriptions overstates the evidence.

\subsection{11.2 What a valid calculation would require}\label{what-a-valid-calculation-would-require}

Let \(C\) be a set of empirically recoded signature clusters rather than prose bullets. A future likelihood ratio would require

\begin{equation}
\log \Lambda
=\sum_{c\in C} w_c
\log
\frac{P(E_c\mid H_D,\mathcal{P})}
{P(E_c\mid H_I,\mathcal{P})},
\tag{14}
\end{equation}

where \(H_D\) and \(H_I\) are declared historical hypotheses, \(\mathcal{P}\) is the named-prior field, and \(w_c\), or a graphical model, accounts for dependence. Calibration would require a sampling frame of comparable papers, blinded recoding, inter-rater reliability, a frozen atom ontology, a declared tail convention, and an access-aware historical model. None presently exists.

The current probability status is therefore \textbf{quantitative-not-calibrated}. The record does not justify a numerical probability of independent or dependent derivation, a Bayes factor, or a sigma claim.

\subsection{11.3 Strongest supported formulation}\label{strongest-supported-formulation}

The absence of a calibrated probability does not erase priority. The strongest supported statement is:

\begin{quote}
Public, time-stamped sources show that Blumberg described the bounded mass-associated microstate -\textgreater{} spatial temporal / statistical gradient -\textgreater{} biased microscopic motion -\textgreater{} gravitational attraction sequence in 2022 and restated it in 2024, before Carney et al.'s February 2025 quantum entropic-force model. The target paper reaches a partially equivalent causal interface through different and more explicit quantum-thermal machinery. That target machinery is a substantive advance, but it does not remove Blumberg's earlier six-node partial priority on the shared sequence.
\end{quote}

This statement is about content and dates. It does not allege copying, access, or intent.

\section{12. What each program contributes}\label{what-each-program-contributes}

\subsection{12.1 QEM contributions preserved}\label{qem-contributions-preserved}

Carney et al.~contribute explicit Hamiltonian models, thermal mediator spectra, local and nonlocal constructions, bath coupling, timescale reduction, Lindblad operators, parameter bounds, and entanglement / noise experiments. Those elements are not retroactively assigned to QGTCD.

\subsection{12.2 QGTCD contributions preserved}\label{qgtcd-contributions-preserved}

QGTCD's distinctive prior contribution in this comparison is the temporal-opportunity causal operator: greater local temporal measure changes directional odds and accumulated microscopic deviations produce motion toward mass. The 2024 local quadrant narrative makes that operation especially clear. QEM does not contain the same time-density ontology, but it later occupies a partially equivalent functional route from a mass-conditioned microscopic statistical state to a gradient-driven Newtonian output.

\subsection{12.3 SIT contribution and margin}\label{sit-contribution-and-margin}

SIT version 1 publicly connected information density, coherence / decoherence, local time density, and gravitational effects on February 10, 2025 {[}15{]}. That is before the February 24 arXiv submission, but the margin is short and the broad information-gravity atom has older named priors. SIT therefore strengthens the continuity of Blumberg's program; the safer priority anchor for the target comparison remains QGTCD's 2022-2024 operator chain.

\section{13. Shared scientific obligations}\label{shared-scientific-obligations}

\subsection{13.1 Source law}\label{source-law}

Both programs must explain why their latent scalar is sourced by ordinary matter in the claimed form. Choosing mediator spectra or mapping a temporal field to a Newtonian potential defines a model; neither alone shows that nature selects it.

\subsection{13.2 Universality}\label{universality}

QGTCD's current drift coefficient contains (D), creating a preparation-dependence problem unless further dynamics remove it. QEM must show how mediator parameters preserve existing universality constraints while leaving a measurable residual. Equivalence-principle tests sharply constrain both possibilities {[}19{]}.

\subsection{13.3 Relativistic completion}\label{relativistic-completion}

An inverse-square force is not a relativistic gravity theory. Light propagation, redshift, covariance, radiation, conservation, strong fields, and post-Newtonian limits remain promotion gates.

\subsection{13.4 Noise}\label{noise}

QEM has the stronger current noise account because it derives jump operators and force fluctuations from a microscopic open system. QGTCD has path variance from a classical kernel, not a quantum-gravity decoherence spectrum. A future synthesis would need to show whether temporal-opportunity dynamics can arise as a controlled limit of a quantum mediator model and what new noise law survives.

\subsection{13.5 Discriminating experiments}\label{discriminating-experiments}

A useful joint experiment program would preregister model-specific likelihoods for:

\begin{enumerate}
\def\labelenumi{\arabic{enumi}.}
\tightlist
\item
  force-noise and heating spectra that constrain QEM mediator-bath parameters;
\item
  matched probe preparations with different diffusion or relaxation properties that test QGTCD implementation dependence;
\item
  entanglement witnesses that exclude specified QEM regions without being overread as proof of one ontology;
\item
  clock or interferometer measurements that distinguish a temporal-opportunity source law from an ordinary thermal mediator.
\end{enumerate}

\section{14. Limitations}\label{limitations}

This paper audits the principal QEM paper, three earlier Carney-lineage gravity-test sources, one related classical-quantum source, Blumberg's public QGTCD genealogy, the first Figshare SIT version, and a private AI-chat comparison index. It does not claim that every publication by every target coauthor has been exhausted. The newest relevant entropic source recovered is the 2025 QEM paper; a 2026 revision of a separate interferometer paper was also checked and remained an EFT-graviton calculation.

The private AI-chat index scanned 5,921 source variants grouped into 3,595 logical families. QEM-related mechanical routing found 35 \texttt{Carney} candidate events across three variants / two logical families, sixteen exact-title events across five variants / three families, fifty-four \texttt{QMEF} events in one family, and thirty-one \texttt{entropic\ forces} events across twelve variants / eight families. These overlapping counts cannot be added. Only exact relevant spans were reviewed, and no private chat body is evidence for public priority. The chats are discovery and audit-history records; public conclusions are grounded in fixed public sources.

The atom count measures bounded content coverage, not independence or statistical confidence. Six Micah-earlier rows share causal dependencies. A future meta-audit may split or merge atoms, but it must preserve this draft and its evidence as an append-only historical judgment rather than silently deleting it.

No new experimental data are reported. The governed development package contains hash-frozen target PDFs and text extractions, a source manifest, an atom ledger, a search census, preserved predecessor drafts, a claim / source / genealogy atlas, and validation scripts. The manuscript and its exact-file package are locally final-preprint ready under the paper-control standard. Independent review remains welcome but is not represented as a completed event. A Zenodo, OSF, DOI, website, or other repository release remains a separate explicit act and is not performed by this manuscript.

\subsection{14.1 Losslessness, authorship, and rights}\label{losslessness-authorship-and-rights}

Draft 7 remains preserved as the immediate source-audited predecessor. This successor retains its four-lens findings, fourteen-atom ledger, nineteen references, target-specific advances, named-prior boundaries, quantitative-not-calibrated status, search census, and scientific failure conditions. The revision repairs malformed inline mathematics, makes the units in the gradient-interface translation explicit, adds the July 27, 2026 newest-source check, and states the local final-preprint status. These changes clarify the existing argument without deleting the earlier favorable or unfavorable findings.

The scientific theories, hypotheses, and historical source materials attributed to Micah Blumberg are his work. Computational tools were used in bounded editorial roles for source routing, mathematical checking, consistency testing, layout, and package validation. Blumberg remains the author and is responsible for the paper's claims and errors.

No third-party figures or article tables are reproduced. The comparison tables are original analytical summaries. Equations from the target literature are either standard definitions or bounded reconstructions presented with citations. Third-party source PDFs are evidence inputs and are excluded from the public release candidate. License selection remains an author / provider decision.

\section{14A. Dated Status Update, Conditional Thermal Calibration, and Secondary Discriminating Observables (November 2026)}\label{a.-dated-status-update-conditional-thermal-calibration-and-secondary-discriminating-observables-november-2026}

\textbf{Dated scientific status (November 20, 2026):} The July 27, 2026 edition was publicly deposited on July 29, 2026 as Paper 062 (Zenodo DOI 10.5281/zenodo.21693943, \url{https://pub.svgn.io/papers/paper-62/}). This update preserves the historical source window and the fundamental distinction between a classical overdamped QGTCD temporal-opportunity model and a quantum mediator-and-bath entropic-force model; a full new physical-reduction paper remains deferred until a microscopic bath-to-QGTCD reduction is derived.

\subsection{Exact conditional thermal calibration}\label{exact-conditional-thermal-calibration}

Within the bounded overdamped QGTCD drift relation \(b = 2\beta D_x \nabla u\), suppose an independently justified thermal bath model supplies Einstein--Smoluchowski mobility relations \(b = \mu F\) (implying physical friction \(\zeta = 1/\mu\)) and \(D_x = \mu k_B T\) together with a constant-gradient driving force \(F = c_A \nabla u\). Matching the mean terminal drift \(b = 2\beta \mu k_B T \nabla u = \mu c_A \nabla u\) forces the exact conditional calibration (verified symbolically in \nolinkurl{verify_paper062_v2.py}):

\begin{equation}
c_A = 2\beta k_B T.
\end{equation}

\subsection{Secondary discriminating observables and physical-regime validity bounds}\label{secondary-discriminating-observables-and-physical-regime-validity-bounds}

Matching a single observed terminal drift \(b_0\) at fixed \(|\nabla u|\) constrains only the product \(\beta \mu k_B T\), leaving bath temperature \(T\) and mobility \(\mu\) degenerate unless secondary observables predicted by the \textbf{same} \((\mu, T, \beta)\) parameters are measured simultaneously. Holding \(b_0\), \(|\nabla u|\), \(\beta\), particle mass \(m\), and separation \(\Delta x\) fixed while varying \(T\) and adjusting \(\mu = b_0 / (2\beta k_B T |\nabla u|)\) inversely keeps the spatial diffusion coefficient \(D_x = \mu k_B T = b_0/(2\beta|\nabla u|)\) identical across models. Under constant-force terminal drift (requiring observation times \(t_{\mathrm{obs}} \gg \tau_p = m/\zeta = m\mu\)), the driven viscous dissipation power is \(P_{\mathrm{diss}} = F \cdot b_0 = \mu F^2 = 2\beta k_B T\, b_0 |\nabla u|\) (scaling \textbf{linearly} in \(T\)).

From \(\zeta = 1/\mu\), the consistent Brownian momentum-diffusion coefficient is \(D_{pp} = \zeta k_B T = k_B T / \mu\). In the high-temperature Markovian Brownian-motion limit, the corresponding decay rate of an \textbf{off-diagonal density-matrix amplitude} \(|\rho(x, x')|\) at spatial separation \(\Delta x = |x - x'|\) is (note that a squared-coherence observable \(|\rho(x, x')|^2\) decays at \(2\Gamma_{\mathrm{dec}}\)):

\begin{equation}
\Gamma_{\mathrm{dec}}(\Delta x) = \frac{D_{pp}}{\hbar^2}(\Delta x)^2 = \frac{k_B T}{\mu \hbar^2}(\Delta x)^2 = \frac{2\beta (k_B T)^2 |\nabla u|}{b_0 \hbar^2}(\Delta x)^2,
\end{equation}

which scales \textbf{quadratically} in \(T\) (\nolinkurl{verify_paper062_v2.py}). Evaluating a representative mesoscopic sphere (\(m = 10^{-15}\,\mathrm{kg}\), \(\Delta x = 10\,\mathrm{nm}\), matched constant-force terminal drift \(b_0 = 10^{-6}\,\mathrm{m/s}\), \(|\nabla u| = 1\,\mathrm{m}^{-1}\), \(\beta = 1\), so \(D_x = 5.00 \times 10^{-7}\,\mathrm{m^2/s}\) is held constant across rows while \(T\) and \(\mu\) vary inversely) yields:

{\def\LTcaptype{none} % do not increment counter
\begingroup\tiny\setlength{\tabcolsep}{2pt}
\begin{longtable}[]{@{}
  >{\raggedleft\arraybackslash}p{(\linewidth - 10\tabcolsep) * \real{0.1667}}
  >{\raggedleft\arraybackslash}p{(\linewidth - 10\tabcolsep) * \real{0.1667}}
  >{\raggedleft\arraybackslash}p{(\linewidth - 10\tabcolsep) * \real{0.1667}}
  >{\raggedleft\arraybackslash}p{(\linewidth - 10\tabcolsep) * \real{0.1667}}
  >{\raggedleft\arraybackslash}p{(\linewidth - 10\tabcolsep) * \real{0.1667}}
  >{\raggedleft\arraybackslash}p{(\linewidth - 10\tabcolsep) * \real{0.1667}}@{}}
\toprule\noalign{}
\begin{minipage}[b]{\linewidth}\raggedleft
\(T\) (\(\mathrm{K}\))
\end{minipage} & \begin{minipage}[b]{\linewidth}\raggedleft
\(\mu\) (\(\mathrm{m/(N\cdot s)}\))
\end{minipage} & \begin{minipage}[b]{\linewidth}\raggedleft
\(\tau_p = m\mu\) (\(\mathrm{s}\))
\end{minipage} & \begin{minipage}[b]{\linewidth}\raggedleft
\(P_{\mathrm{diss}}\) (\(\mathrm{W}\))
\end{minipage} & \begin{minipage}[b]{\linewidth}\raggedleft
Amplitude \(\Gamma_{\mathrm{dec}}\) (\(\mathrm{s}^{-1}\))
\end{minipage} & \begin{minipage}[b]{\linewidth}\raggedleft
\(\Gamma_{\mathrm{dec}} / (k_B T/\hbar)\)
\end{minipage} \\
\midrule\noalign{}
\endhead
\bottomrule\noalign{}
\endlastfoot
\(0.01\) & \(3.62 \times 10^{18}\) & \(3{,}622\) & \(2.76 \times 10^{-31}\) & \(3.43 \times 10^{8}\) & \(0.262\) \\
\(1.00\) & \(3.62 \times 10^{16}\) & \(36.2\) & \(2.76 \times 10^{-29}\) & \(3.43 \times 10^{12}\) & \(26.2\) \\
\(300.0\) & \(1.21 \times 10^{14}\) & \(0.121\) & \(8.28 \times 10^{-27}\) & \(3.09 \times 10^{17}\) & \(7{,}855\) \\
\end{longtable}
\endgroup
}

\textbf{Critical physical-regime qualification:} While the formal parameter scalings separate the \(10\,\mathrm{mK}\) and \(300\,\mathrm{K}\) models by \(3.00 \times 10^4\times\) in driven terminal-drift power and \(9.00 \times 10^8\times\) in nominal decoherence rate, the table exposes two essential regime constraints that any concrete experiment must satisfy:

\begin{enumerate}
\def\labelenumi{\arabic{enumi}.}
\tightlist
\item
  \textbf{Breakdown of naive short-time Markovian decoherence:} At \(1\,\mathrm{K}\) and \(300\,\mathrm{K}\), the nominal amplitude rate exceeds the thermal bath rate \(\omega_{\mathrm{th}} = k_B T/\hbar\) by factors of \(26.2\) and \textbf{\(7{,}855\)}, respectively (and reaches \(0.262\,\omega_{\mathrm{th}}\) even at \(10\,\mathrm{mK}\)). When \(\Gamma_{\mathrm{dec}} \gtrsim k_B T/\hbar\), the nominal decay time \(1/\Gamma_{\mathrm{dec}}\) is shorter than the thermal memory time \(\hbar/(k_B T)\), so the Markovian formula \textbf{cannot} be applied blindly to initial decay without specifying the bath spectral density, high-frequency cutoff \(\Omega_c\), memory time, and measurement interval \(\Delta t\).
\item
  \textbf{Terminal-drift equilibration protocol:} Across the three rows, the momentum relaxation time \(\tau_p = m\mu\) spans nearly four orders of magnitude, from \(0.121\,\mathrm{s}\) at \(300\,\mathrm{K}\) up to \(3{,}622\,\mathrm{s}\) (\(\approx 1\,\mathrm{hr}\)) at \(10\,\mathrm{mK}\). Measuring constant-force terminal drift \(b_0 = \mu F\) and driven power \(P_{\mathrm{diss}} = F b_0\) requires an observation protocol with \(t_{\mathrm{obs}} \gg \tau_p\) (or an explicitly underdamped transient model).
\end{enumerate}

A genuinely discriminating physical reduction must therefore specify bath cutoff, memory time, and equilibration timescale alongside predicting noise, heating, or decoherence from the same parameters that explain mean drift.

\section{15. Conclusion}\label{conclusion}

The corrected comparison is neither full identity nor full separation. Carney et al.'s quantum entropic-force models and Blumberg's QGTCD differ in state space, internal dynamics, mathematical maturity, and empirical specificity. Those differences matter. They do not erase a bounded shared causal architecture.

Blumberg's public 2022-2024 record already contains six partial-priority nodes spanning a mass-associated microscopic substrate, spatially unequal time or statistical measure, a gradient-controlled directional bias, accumulation into motion toward mass, an emergent non-graviton route, and a local discretized implementation. Carney et al.'s 2025 paper later realizes a partially equivalent microstate -\textgreater{} statistical scalar -\textgreater{} gradient -\textgreater{} Newtonian sequence with qubits or oscillators and adds explicit quantum-thermal dynamics, noise, and experiments.

All four lenses support the bounded match. The mechanism lens passes at the mean scalar-gradient interface. The narrative lens passes on the four-stage causal sequence. The conclusion-first lens recovers the same shared operations while preserving two-way residuals. The longitudinal lens shows that the entropic free-energy sequence is absent from the checked 2021-2024 Carney / Taylor stages and appears in 2025 after Blumberg's public QGTCD sources.

The priority result is therefore clear at the stated granularity: \textbf{Micah Blumberg is earlier on the six-node QGTCD partial causal sequence; Carney et al.~are later on that shared interface and contribute a more explicit quantum mediator-bath construction with substantial target-specific advances.} The numerical probability of historical dependence is not identified by the present evidence, and no allegation of copying or intent is made.

\section{References}\label{references}

\begin{enumerate}
\def\labelenumi{\arabic{enumi}.}
\tightlist
\item
  D. Carney, M. Karydas, T. Scharnhorst, R. Singh, and J. M. Taylor, ``On the Quantum Mechanics of Entropic Forces,'' \emph{Physical Review X} \textbf{15}, 031038 (2025), \href{https://doi.org/10.1103/PhysRevX.15.031038}{doi:\nolinkurl{10.1103/PhysRevX.15.031038}}; arXiv:2502.17575, \url{https://arxiv.org/abs/2502.17575}.
\item
  D. Carney, H. Muller, and J. M. Taylor, ``Using an Atom Interferometer to Infer Gravitational Entanglement Generation,'' \emph{PRX Quantum} \textbf{2}, 030330 (2021), \href{https://doi.org/10.1103/PRXQuantum.2.030330}{doi:\nolinkurl{10.1103/PRXQuantum.2.030330}}.
\item
  D. Carney, V. Domcke, and N. L. Rodd, ``Graviton Detection and the Quantization of Gravity,'' \emph{Physical Review D} \textbf{109}, 044009 (2024), arXiv:2308.12988, \url{https://arxiv.org/abs/2308.12988}.
\item
  D. Carney, M. Karydas, and A. Sivaramakrishnan, ``Response of Interferometers to the Vacuum of Quantum Gravity,'' \emph{Physical Review D} \textbf{113}, 106002 (2026), arXiv:2409.03894, \url{https://arxiv.org/abs/2409.03894}.
\item
  D. Carney and A. Matsumura, ``Classical-Quantum Scattering,'' \emph{Classical and Quantum Gravity} \textbf{42}, 135010 (2025), arXiv:2412.04839, \url{https://arxiv.org/abs/2412.04839}.
\item
  T. Jacobson, ``Thermodynamics of Spacetime: The Einstein Equation of State,'' \emph{Physical Review Letters} \textbf{75}, 1260-1263 (1995), \href{https://doi.org/10.1103/PhysRevLett.75.1260}{doi:\nolinkurl{10.1103/PhysRevLett.75.1260}}.
\item
  E. Verlinde, ``On the Origin of Gravity and the Laws of Newton,'' \emph{Journal of High Energy Physics} \textbf{2011}, 29 (2011), \href{https://doi.org/10.1007/JHEP04(2011)029}{doi:\nolinkurl{10.1007/JHEP04(2011)029}}.
\item
  D. Kafri, J. M. Taylor, and G. J. Milburn, ``A Classical Channel Model for Gravitational Decoherence,'' \emph{New Journal of Physics} \textbf{16}, 065020 (2014), \href{https://doi.org/10.1088/1367-2630/16/6/065020}{doi:\nolinkurl{10.1088/1367-2630/16/6/065020}}.
\item
  M. Blumberg, ``QGTCD Source and Evidence Map,'' governed source map (2026), preserving the public Git genealogy and its immutable commit locators.
\item
  M. Blumberg, \texttt{a0253z.md}, public Git state of July 26, 2022, commit \nolinkurl{92dc0fd1aad256664df035711b58e7d8e273f327}, \url{https://github.com/v5ma/selfawarenetworks/blob/92dc0fd1aad256664df035711b58e7d8e273f327/a0253z.md}.
\item
  M. Blumberg, \texttt{a0253z.md}, public Git state of July 30, 2022, commit \nolinkurl{c4e94e9fcd049c97d47ef6a5f2f8c29485f5a98d}, \url{https://github.com/v5ma/selfawarenetworks/blob/c4e94e9fcd049c97d47ef6a5f2f8c29485f5a98d/a0253z.md}.
\item
  M. Blumberg, ``New Unified Field Theory: Quantum Gradient Time Crystal Dilation,'' \emph{SVGN.io}, January 29, 2024, \url{https://www.svgn.io/p/a-new-unified-field-theory-called}.
\item
  M. Blumberg, ``Quantum Gradient Time Crystal Dilation Breaks the Assumption that Gravity Equals Metric Curvature Alone,'' \emph{SVGN.io}, February 15, 2024, \url{https://www.svgn.io/p/quantum-gradient-time-crystal-dilation}.
\item
  M. Blumberg, ``Explain It to Me Like I Am Six: Quantum Gradient Time Crystal Dilation Theory, Part 3,'' \emph{SVGN.io}, March 28, 2024, \url{https://www.svgn.io/p/explain-it-to-me-like-i-am-six-quantum}.
\item
  M. Blumberg, ``Super Information Theory,'' Figshare version 1, posted February 10, 2025, \url{https://figshare.com/articles/journal_contribution/Super_Information_Theory/28379318/1}.
\item
  M. Blumberg, ``Quantum Gradient Time Crystal Dilation: A Minimal Temporal-Measure-Biased Random Walk and Its Drift-Diffusion Generator,'' private source-bounded manuscript (2026); archival lineage \href{https://doi.org/10.5281/zenodo.21459858}{doi:\nolinkurl{10.5281/zenodo.21459858}}.
\item
  S. Bose et al., ``Spin Entanglement Witness for Quantum Gravity,'' \emph{Physical Review Letters} \textbf{119}, 240401 (2017), \href{https://doi.org/10.1103/PhysRevLett.119.240401}{doi:\nolinkurl{10.1103/PhysRevLett.119.240401}}.
\item
  C. Marletto and V. Vedral, ``Gravitationally Induced Entanglement between Two Massive Particles Is Sufficient Evidence of Quantum Effects in Gravity,'' \emph{Physical Review Letters} \textbf{119}, 240402 (2017), \href{https://doi.org/10.1103/PhysRevLett.119.240402}{doi:\nolinkurl{10.1103/PhysRevLett.119.240402}}.
\item
  P. Touboul et al., ``MICROSCOPE Mission: Final Results of the Test of the Equivalence Principle,'' \emph{Physical Review Letters} \textbf{129}, 121102 (2022), \href{https://doi.org/10.1103/PhysRevLett.129.121102}{doi:\nolinkurl{10.1103/PhysRevLett.129.121102}}.
\end{enumerate}

\end{document}
